\section{AI 소스 코드별 API Reference}

본 부록은 \code{AI/} 디렉터리의 소스(\code{ai.py})를 정리한다. 브라우저에서
동작하는 문법 기반 자연어 파서(\code{ai\_helper.js},
\code{ai\_helper\_grammar.js})는 \code{Web/}에 위치하므로 부록 A2에 수록한다.

\subsection{ai.py}

\noindent\textbf{역할}: EXAONE-3.5-2.4B 로컬 LLM을 래핑하여, 초등학생의 블록
코딩\textbf{·}로봇 조작 질문에 한국어로 답하는 코드 도우미 PoC.

\noindent\textbf{모델\textbf{·}생성 설정}
\begin{itemize}
  \item \code{model\_name = "LGAI-EXAONE/EXAONE-3.5-2.4B-Instruct"}
  \item 로딩: \code{dtype=torch.bfloat16}, \code{trust\_remote\_code=True},
        \code{device\_map="auto"}
  \item 스트리밍 생성: \code{max\_new\_tokens=2048}, \code{temperature=0.6},
        \code{top\_p=0.95}, \code{do\_sample=True}
  \item 비스트리밍 생성: \code{max\_new\_tokens=256}, \code{temperature=0.6},
        \code{top\_p=0.95}, \code{do\_sample=True}
\end{itemize}

\noindent\textbf{class CodeAssistant} --- 어린이 대상 블록 코딩 질의응답 도우미.
속성: \code{self.model}, \code{self.tokenizer}, \code{self.device}.

\begin{longtable}{S{8.6cm} L{5.4cm}}
\toprule
\rowcolor{tablehead}\normalfont\sffamily\bfseries 시그니처 & \normalfont\sffamily\bfseries 설명 \\
\midrule
\endhead
def \_\_init\_\_(self, model, tokenizer, device="cuda"): & 모델\textbf{·}토크나이저\textbf{·}디바이스 설정 \\
def ask\_about\_code(self, user\_query): & 질문 처리\textbf{·}응답 생성(스트리밍/비스트리밍) \\
\bottomrule
\end{longtable}

\noindent\textbf{모듈 함수}: \code{def call\_ai(prompt):} ---
\code{CodeAssistant}를 인스턴스화하여 질의를 위임하는 편의 함수.

\noindent\textbf{시스템 프롬프트}: 모델이 초등학생 대상 라즈베리파이 로봇 조작
튜터로서, Google Blockly 블록 코딩을 쉬운 한국어로 설명하도록 지시한다.
로봇 제어 블록(램프\textbf{·}메시지\textbf{·}상태\textbf{·}거리/자기 센서\textbf{·}부저\textbf{·}이동\textbf{·}
정지)과 제어 블록(대기\textbf{·}반복\textbf{·}변수\textbf{·}기초 연산)을 안내한다.

\noindent\textbf{외부 의존성}: \code{torch}(\code{bfloat16}),
\code{transformers.AutoModelForCausalLM}, \code{transformers.AutoTokenizer},
\code{transformers.TextIteratorStreamer}, \code{threading.Thread}(스트리밍 생성),
\code{pathlib.Path}.

\vspace{2em}
{\color{rulegray}\hrule height 0.6pt}
\vspace{0.5em}
\begin{flushright}
\small\sffamily\color{rulegray}
--- 부록 끝 ---
\end{flushright}
